%! Skills Focus: Selecting, Implementing, and Communicating Inference Procedures — Unit 7, Lesson 7.10 — Guided Notes (Answer Key)
\documentclass[10pt]{article}
\usepackage{apstats-article}
\usepackage{apstats-key}   % pulls in -boxes; do NOT also load -boxes
\newcommand{\TallMath}[1]{$\displaystyle #1\rule[-1.4em]{0pt}{3.2em}$}

\begin{document}

\pageheader{Unit 7, Lesson 7.10}{Guided Notes --- Answer Key}
\namedateperiod

\vspace{0.12in}

\begin{objectivebox}
By the end of this lesson, I will be able to:
\begin{itemize}
  \item \ansline{Select the correct inference procedure by identifying the parameter type and number of samples.}
  \item \ansline{Verify the appropriate conditions and communicate whether each is met using numbers from the context.}
  \item \ansline{Write a complete four-step inference procedure (State, Plan, Do, Conclude) using AP-exam language.}
\end{itemize}
\end{objectivebox}

\begin{vocabbox}
\begin{description}
  \item[\termblanklong{Inference procedure selection}]
        To choose correctly: (1) identify the \ans{parameter} (proportion or mean); (2) count
        the \ans{samples} (one or two); (3) decide whether to \ans{estimate} or \ans{test}.
  \item[\termblanklong{Matched-pairs $t$-test}]
        Used when data come in \ans{pairs} (same subject measured twice, or two paired
        subjects). Compute \ans{differences $d_i$}, then run a \ans{one-sample} $t$-test on the differences.
        The parameter is \ans{$\mu_d$}.
\end{description}
\end{vocabbox}

\begin{hookbox}
A sports scientist records resting heart rates for 20 elite cyclists and 20 recreational
cyclists. She wants to know whether elite cyclists have a lower \emph{mean} resting heart rate.

What procedure should she use? \ans{two-sample $t$-test for $\mu_1 - \mu_2$}

How did you decide? Two questions:
\begin{enumerate}
  \item Parameter type: \ans{mean (quantitative data, not a proportion)}
  \item Number of samples: \ans{two independent groups (elite vs.\ recreational)}
\end{enumerate}
\end{hookbox}

\begin{notesbox}{The Procedure Selection Framework}
\textbf{Ask these questions in order:}

\vspace{0.06in}
\renewcommand{\arraystretch}{2.0}
\begin{tabularx}{\linewidth}{@{} p{0.28\linewidth} p{0.20\linewidth} p{0.16\linewidth} X @{}}
  \rowcolor{sky}
  \textbf{Parameter} & \textbf{\# Samples} & \textbf{Goal} & \textbf{Procedure} \\ \hline
  Proportion $p$ & One & CI & \ans{one-sample $z$-interval for $p$} \\
  Proportion $p$ & One & Test & \ans{one-sample $z$-test for $p$} \\
  $p_1-p_2$ & Two & CI & \ans{two-sample $z$-interval for $p_1-p_2$} \\
  $p_1-p_2$ & Two & Test & \ans{two-sample $z$-test for $p_1-p_2$} \\
  Mean $\mu$ & One & CI & \ans{one-sample $t$-interval for $\mu$} \\
  Mean $\mu$ & One & Test & \ans{one-sample $t$-test for $\mu$} \\
  $\mu_1-\mu_2$ & Two & CI & \ans{two-sample $t$-interval for $\mu_1-\mu_2$} \\
  $\mu_1-\mu_2$ & Two & Test & \ans{two-sample $t$-test for $\mu_1-\mu_2$} \\
\end{tabularx}
\end{notesbox}

\begin{notesbox}{The Four-Step AP Write-Up}
\renewcommand{\arraystretch}{1.8}
\begin{tabularx}{\linewidth}{@{} p{0.10\linewidth} p{0.22\linewidth} X @{}}
  \rowcolor{sky}
  \textbf{Step} & \textbf{Name} & \textbf{What to write} \\ \hline
  1 & \ans{State} & Name the procedure; state $H_0$ and $H_a$ (or parameter for a CI); define symbols. \\
  2 & \ans{Plan} & List each condition; verify with \ans{numbers from context}; state whether met. \\
  3 & \ans{Do} & Show formula + substitution; report \ans{$p$-value} (or interval). \\
  4 & \ans{Conclude} & ``Because $p$-value \ans{$<$} $\alpha$, I \ans{reject} $H_0$.
                     There \ans{is} sufficient evidence that \ans{[restate $H_a$ in context]}.'' \\
\end{tabularx}
\end{notesbox}

\begin{practicebox}
\textbf{Partner Practice.} \;
A school district randomly selects 40 elementary schools and 40 middle schools and records
the mean reading-score gain (points) for each school over one year. An analyst wants to know
whether elementary schools showed a larger mean gain than middle schools.

\begin{enumerate}
  \item Parameter: \ans{$\mu_1 - \mu_2$ (difference of means)} \quad \# samples: \ans{two} \quad Goal: \ans{test}
  \item Procedure name: \ans{two-sample $t$-test for $\mu_1 - \mu_2$}
  \item $H_0$: \ans{$\mu_1 - \mu_2 = 0$} \quad $H_a$: \ans{$\mu_1 - \mu_2 > 0$}
  \item List the three conditions and state how each is verified:
        \begin{itemize}
          \item Random: \ansline{Both groups were randomly selected --- Random condition met.}
          \item 10\%: \ansline{$40 \le 10\%$ of all elementary (and middle) schools in the district (assuming large population) --- met.}
          \item Normal: \ansline{$n_1 = n_2 = 40 \ge 30$ --- CLT applies; Normal condition met.}
        \end{itemize}
\end{enumerate}
\end{practicebox}

\end{document}
